\chapter{Arrancar el hub y gobernar un \emph{namespace}}
\label{ch:hub}

\cod{alexandria-server} es el hub. Tiene dos caras: una API de gestión autenticada bajo \cod{/v1} y un árbol de
publicación público e inmutable bajo \cod{/pub}. Este capítulo lo arranca, crea el \emph{namespace} de Flexgrid
y le asigna una \emph{steward}.

\section{El laboratorio}

Todo ocurre en \cod{tutorial/manual/examples/lab}. El fichero \cod{env.sh} define las variables y unas pocas
funciones que envuelven a \cod{curl}; ninguna esconde una llamada.

\begin{lstlisting}[style=sh]
cd tutorial/manual/examples/lab
source env.sh
\end{lstlisting}

\noindent\begin{tabularx}{\linewidth}{@{}L{3.3cm}L{5.1cm}Y@{}}
\toprule
\textbf{Variable} & \textbf{Valor} & \textbf{Qué es} \\
\midrule
\texttt{\$BIN} & \cod{../../../../bin} & Los binarios compilados. \\
\texttt{\$HUB} & \cod{https://127.0.0.1:18443} & El hub. \\
\texttt{\$IDP} & \cod{http://127.0.0.1:19400} & El emisor de tokens. \\
\texttt{\$CONNECTOR} & \cod{http://127.0.0.1:18282} & El conector de mentira. \\
\texttt{\$INGRESS} & \cod{http://127.0.0.1:19080} & La entrada de tráfico del cliente. \\
\texttt{\$ADMINAPI} & \cod{http://127.0.0.1:19090} & La API de administración del cliente. \\
\bottomrule
\end{tabularx}

\subsection{Un reloj que se puede mover}

El ciclo de vida de una versión dura semanas. Para recorrerlo en minutos, todos los procesos aceptan
\cod{-clock-file}: leen la hora de un fichero en lugar del reloj del sistema. La función \cod{clock} escribe en
él.

\begin{lstlisting}[style=sh]
clock 2026-09-01T09:00:00Z
\end{lstlisting}

\begin{trampaHabitual}[Todos con el mismo reloj]
El hub, el cliente y el emisor de tokens tienen que leer el \textbf{mismo} fichero. Un token emitido con la hora
real del sistema está, para un hub que cree vivir en septiembre, caducado o aún no vigente, y la respuesta es
\cod{401}. Por eso \cod{tok} pasa también \cod{--clock-file}.
\end{trampaHabitual}

\section{La configuración del hub}

\cod{lab/server.yaml} usa la disposición por secciones. Lo que lo distingue de \cod{config/server.yaml} son tres
cosas: los puertos, el compilador y el directorio de copias vendorizadas. Aquí se muestra condensado.

\begin{lstlisting}[style=yaml]
common_config:
  hosts:
    http: { protocol: "https", url: "127.0.0.1", port: "18443", internal_port: "18443" }
  db:
    db_type: "SQLite"
    database: "file:var/hub/hub.db"
    blobs: "file:var/hub/blobs"
  connection:
    tls: { cert: "file:var/tls/hub.crt", key: "file:var/tls/hub.key" }

auth_config:
  enabled: true
  issuer: "http://127.0.0.1:19400"
  audience: "alexandria-vhub"
  roles_claim: "roles"
  admin_role: "vhub-admin"

did_config:
  type: "key"
  signing_key: "file:var/hub/hub.key"

vhub_config:
  rdf_context: "self-hosted"
  compiler: "dsl"
  vendor_dir: "file:../vendor"
\end{lstlisting}

\noindent\begin{tabularx}{\linewidth}{@{}L{5.3cm}Y@{}}
\toprule
\textbf{Campo} & \textbf{Significado} \\
\midrule
\cod{common_config.hosts.http} & Dirección pública y puerto. De aquí sale la URL base, prefijo de toda IRI que emite el hub. \\
\cod{common_config.db} & Base de datos SQLite y directorio de blobs direccionados por contenido. \\
\cod{connection.tls} & Certificado y clave. Si no existen, se genera un certificado de laboratorio. \\
\cod{auth_config} & El hub es un servidor de recursos OAuth 2.0: emisor aceptado, audiencia y la reclamación que marca al administrador. \\
\cod{did_config.signing_key} & Clave Ed25519 con la que se firma. Si no existe, se genera. El DID del hub es por defecto el \cod{did:key} de esa clave. \\
\cod{vhub_config.compiler} & \cod{poc}, por defecto: borradores JSON. \cod{dsl}: borradores YAML del lenguaje de perfiles. \\
\cod{vhub_config.vendor_dir} & Copias locales para los \emph{bindings} vendorizados. \\
\cod{vhub_config.rdf_context} & \cod{self-hosted}: el hub sirve su propio contexto JSON-LD en \cod{/pub/_vhub/v1/context.jsonld}. \\
\bottomrule
\end{tabularx}

\medskip
Las rutas \cod{file:} se resuelven desde el directorio del fichero de configuración. El documento se valida al
arrancar contra \cod{deploy/server.schema.json}; con un error, el hub no arranca.

\begin{notaPractica}[Dos disposiciones, el mismo contenido]
Los dos binarios aceptan YAML o JSON, y dos disposiciones equivalentes: la de secciones, que usa este tutorial,
y la plana de la especificación (\cod{deploy/server.json}), donde los mismos campos se llaman \cod{listen},
\cod{base}, \cod{signing.key}, \cod{auth.issuers} o \cod{storage.database}. En la plana, \cod{compiler} y
\cod{vendor_dir} van en la raíz.
\end{notaPractica}

\section{Arrancar}

\begin{lstlisting}[style=sh]
$BIN/fake-idp serve --addr 127.0.0.1:19400 --issuer $IDP --key idp.key --clock-file clock > idp.log 2>&1 &
$BIN/fake-connector -listen 127.0.0.1:18282 > connector.log 2>&1 &
$BIN/alexandria-server -config server.yaml -clock-file clock > server.log 2>&1 &
\end{lstlisting}

En el primer arranque, el hub crea lo que le falta y lo dice:

\begin{lstlisting}[style=out]
time=2026-09-01T09:00:00.000Z level=INFO msg="generated lab TLS certificate" … cert_path=var/tls/hub.crt note="PoC: not for production"
time=2026-09-01T09:00:00.000Z level=INFO msg="signing key verified in hub DID document" … hub_did=did:key:z6Mkkfgq…
time=2026-09-01T09:00:00.000Z level=INFO msg="server listening" … listen=127.0.0.1:18443 base=https://127.0.0.1:18443 tls_enabled=true
\end{lstlisting}

Tres cosas a notar. Las marcas de tiempo son las del reloj de fichero. El hub comprueba que su clave de firma
es un método de aserción de su propio documento DID antes de aceptar tráfico: nunca firma con una clave que no
haya publicado. Y el \cod{hub_did} hay que apuntarlo, porque el cliente solo aceptará credenciales de ese DID.

\begin{lstlisting}[style=sh]
hubdid
\end{lstlisting}
\begin{lstlisting}[style=out]
did:key:z6MkkfgqPKJK19zucGBNfLJQHUFwEA3mc9skvYsPDbkoNP8e
\end{lstlisting}

El tuyo será distinto: la clave se genera en cada laboratorio nuevo.

\section{Quién puede hacer qué}

\subsection{Tokens}

\texttt{fake-idp token} firma un token de acceso. Con \cod{--admin} añade el rol de administrador del hub.

\begin{lstlisting}[style=sh]
tok --sub ana --admin > admin.token
tok --sub lucia > steward.token
\end{lstlisting}

Ana es la administradora del hub. Lucía, de la distribuidora, será la \emph{steward} de Flexgrid. Por ahora su
token solo dice quién es; el hub decidirá qué puede hacer.

El hub valida cada token en cinco pasos, y el primero que falla responde \cod{401}:

\noindent\begin{tabularx}{\linewidth}{@{}C{1.1cm}Y@{}}
\toprule
\textbf{Paso} & \textbf{Comprobación} \\
\midrule
A1 & Es un JWT de tipo \cod{at+jwt} firmado con un algoritmo asimétrico permitido. \\
A2 & Su emisor es uno de los configurados. \\
A3 & La firma verifica contra las claves públicas del emisor, que se descubren por OIDC. \\
A4 & La audiencia incluye la del hub. \\
A5 & Está dentro de su periodo de validez, con la tolerancia de reloj. \\
\bottomrule
\end{tabularx}

\subsection{La tabla de autorización}

\noindent\begin{tabularx}{\linewidth}{@{}Y C{1.6cm} C{2.3cm} C{1.5cm} C{1.2cm}@{}}
\toprule
\textbf{Operación} & \textbf{Anónimo} & \textbf{Autenticado} & \textbf{Steward} & \textbf{Admin} \\
\midrule
Leer \emph{namespaces}, vocabularios, listado y registros; todo \cod{/pub} & \si & \si & \si & \si \\
Crear un \emph{namespace}; asignar o quitar \emph{stewards}; leer el registro de eventos global & 401 & 403 & 403 & \si \\
Listar los \emph{stewards} de un \emph{namespace} & 401 & 403 & \si & \si \\
Vocabularios, borradores, candidatos, aprobación, revocación y eventos de un \emph{namespace} & 401 & 403 & \si & 403 \\
\bottomrule
\end{tabularx}

\begin{trampaHabitual}[El administrador no gobierna el vocabulario]
La última fila sorprende: el administrador recibe \cod{403} en las operaciones de un \emph{namespace}. Ser
administrador del hub no da poder sobre el contenido de un espacio. Si Ana quisiera publicar versiones de
Flexgrid, tendría que asignarse como \emph{steward}.
\end{trampaHabitual}

Las respuestas reales de los dos casos de error:

\begin{lstlisting}[style=sh]
curl -sk -X POST $HUB/v1/namespaces -d '{"slug":"x"}'
\end{lstlisting}
\begin{lstlisting}[style=out]
{"@context":"https://127.0.0.1:18443/pub/_vhub/v1/context.jsonld",
 "type":"https://w3id.org/alexandria-vhub/problems#API_UNAUTHENTICATED","title":"API_UNAUTHENTICATED",
 "status":401,"detail":"this operation requires a bearer token","code":"API_UNAUTHENTICATED","severity":"error"}
\end{lstlisting}

Un administrador intentando crear un borrador devuelve la misma forma con \cod{API_FORBIDDEN_ROLE} y \cod{403}.
Toda respuesta de \cod{/v1}, errores incluidos, es JSON-LD compacto con un \cod{code} del catálogo cerrado.

\section{Crear el \emph{namespace} y su \emph{steward}}

\begin{lstlisting}[style=sh]
curl -sk -X POST $HUB/v1/namespaces -H "Authorization: Bearer $(cat admin.token)" \
  -d '{"slug":"flexgrid","title":"Flexgrid Energy Data Space"}'
\end{lstlisting}
\begin{lstlisting}[style=out]
{"@context":"https://127.0.0.1:18443/pub/_vhub/v1/context.jsonld","created_at":"2026-09-01T09:00:00Z",
 "id":"https://127.0.0.1:18443/v1/namespaces/flexgrid","slug":"flexgrid",
 "title":"Flexgrid Energy Data Space","type":"Namespace"}
\end{lstlisting}

El \cod{slug} es único e inmutable y forma parte de todas las URL. \cod{_vhub} está reservado.

\begin{lstlisting}[style=sh]
curl -sk -X POST $HUB/v1/namespaces/flexgrid/stewards -H "Authorization: Bearer $(cat admin.token)" \
  -d '{"iss":"http://127.0.0.1:19400","sub":"lucia","label":"Lucía (DSO)"}'
\end{lstlisting}
\begin{lstlisting}[style=out]
{"@context":"…","assigned_at":"2026-09-01T09:00:00Z","assigned_by":"http://127.0.0.1:19400#ana",
 "id":"https://127.0.0.1:18443/v1/namespaces/flexgrid/stewards/ed6edbdc291414db",
 "iss":"http://127.0.0.1:19400","label":"Lucía (DSO)","sub":"lucia","type":"Steward"}
\end{lstlisting}

Una \emph{steward} se identifica por el par emisor y sujeto, no por un nombre. La pertenencia se evalúa en cada
petición: el token de Lucía no ha cambiado, pero desde este instante el hub le deja operar Flexgrid.

\section{El registro de vocabularios}

Además de versiones, un \emph{namespace} puede alojar documentos sueltos con revisiones inmutables: una lista de
códigos, un contexto de terceros fijado, una ontología auxiliar. Hay dos tipos.

\noindent\begin{tabularx}{\linewidth}{@{}L{2cm}Y@{}}
\toprule
\textbf{Tipo} & \textbf{Comportamiento} \\
\midrule
\cod{hosted} & El contenido se sube. Cada subida distinta crea una revisión nueva; subir los mismos bytes no crea nada. \\
\cod{proxy} & El hub descarga un documento externo al crear el vocabulario y al pedir \cod{:refresh}. Nunca contacta con el origen al servir \cod{/pub}. \\
\bottomrule
\end{tabularx}

\medskip
Flexgrid aloja su esquema de regiones, \cod{examples/dsl/regions.ttl}:

\begin{lstlisting}[style=sh]
S=$(cat steward.token)
curl -sk -X POST $HUB/v1/namespaces/flexgrid/vocabularies -H "Authorization: Bearer $S" \
  -d '{"name":"regions","kind":"hosted","media_type":"text/turtle"}'

curl -sk -i -X PUT $HUB/v1/namespaces/flexgrid/vocabularies/regions/content \
  -H "Authorization: Bearer $S" -H 'Content-Type: text/turtle' --data-binary @../dsl/regions.ttl
\end{lstlisting}
\begin{lstlisting}[style=out]
HTTP/2 201
{"@context":"…","created_at":"…","created_by":"http://127.0.0.1:19400#lucia",
 "digest":"ni:///sha-256;zmzyAQLYqxTjnY6aXPy_HRSkqmgxrO7MhJsJkU6ohOI",
 "id":"https://127.0.0.1:18443/pub/flexgrid/vocab/regions/r/1","n":1,"size":432,"type":"Revision"}
\end{lstlisting}

Repetir la misma subida devuelve \cod{200} y no crea revisión. Subir algo que no es Turtle se rechaza:

\begin{lstlisting}[style=out]
{"…","title":"VOCAB_UNPARSEABLE","status":422,"detail":"content is not valid text/turtle","code":"VOCAB_UNPARSEABLE","severity":"error"}
\end{lstlisting}

La revisión es pública de inmediato. La URL sin revisión redirige a la última; la URL con revisión es inmutable:

\begin{lstlisting}[style=sh]
curl -sk -i $HUB/pub/flexgrid/vocab/regions
curl -sk -i $HUB/pub/flexgrid/vocab/regions/r/1
\end{lstlisting}
\begin{lstlisting}[style=out]
HTTP/2 303
cache-control: no-cache
location: /pub/flexgrid/vocab/regions/r/1

HTTP/2 200
cache-control: public, max-age=31536000, immutable
content-type: text/turtle
etag: "ce6cf20102d8ab14e39d8e9a5cfcbf1d14a4aa6831aceecc849b09914ea884e2"
\end{lstlisting}

\begin{ideaClave}
Este patrón se repite en todo el hub: una URL \textbf{estable} que redirige y no se cachea, y una URL
\textbf{inmutable} que se cachea para siempre y cuyo \cod{ETag} es el SHA-256 del contenido. Quien quiera «lo
último» usa la primera; quien necesite reproducibilidad cita la segunda.
\end{ideaClave}

\begin{retos}
\begin{enumerate}[leftmargin=*, itemsep=2pt]
  \item Emite un token para un tercero, \texttt{tok -{}-sub eve}, e intenta crear un vocabulario. ¿Qué código devuelve?
  \item Sube una segunda versión de \cod{regions.ttl} con una región más y comprueba que
        \cod{/pub/flexgrid/vocab/regions} redirige ahora a \cod{r/2} mientras que \cod{r/1} sigue igual.
  \item Intenta crear el \emph{namespace} \cod{flexgrid} otra vez, y luego uno llamado \cod{Flexgrid!}.
\end{enumerate}
\end{retos}
